\chapter*{Abstract}
\addcontentsline{toc}{chapter}{Abstract}

A neural network used as a variational wavefunction is judged by its energy. When that energy
stalls, or two networks reach it by different routes, the number cannot say why: the
representation may be wrong for the problem, the optimiser may be failing, the samples that
define the loss may no longer describe the state, or the physics may have become harder. This
thesis tries to tell those explanations apart in the two-dimensional harmonic quantum dot, for
$N\in\{2,6,12,20\}$ electrons, from the compact droplet at $\omega=1$ to the Wigner molecule at
$\omega=10^{-3}$.

The ansatz is a Slater determinant dressed by a neural Jastrow correlator, which reshapes the
amplitude, and a coordinate backflow, which reshapes the nodes. Trained by residual
pretraining and stochastic reconfiguration, it agrees to within $0.05\%$ with published
reference energies at the confinements where we compare with them, $\omega=1$, $0.5$ and
$0.1$; below that, for $N\ge6$, its energies are predictions.

The correlator can be \emph{separable}, encoding each particle and each pair independently
before pooling them, or \emph{message-passing}, letting the particles exchange information
first. The two reach the same energy, but not the same state. The message-passing correlator
carries a local-energy variance lower by a factor of $1.2$--$3.4$ at every confinement, and
its leading tangent directions lie almost entirely within a handful of physical collective
modes (breathing, quadrupole, correlation hole), where a separable network of equal or
greater size leaves close to half of its tangent weight outside them. Stochastic
reconfiguration and Adam agree to four decimals on models whose metric has condition numbers
of $10^{7}$--$10^{10}$; the failures suggest that natural-gradient preconditioning pays when
the metric can be estimated from a batch, not merely when it is ill-conditioned.

On its own, the strong-form residual that a physics-informed network would minimise could not
train the backflow in any variant we tried: its learning signal lies at the nodes, where the
local energy diverges, and it stalls near $0.3\%$. A weak form, the variational energy
gradient estimated by importance sampling from an analytic proposal, keeps the Markov chain
out of the gradient and has no such floor. With a proposal refitted during training it trains
the six-electron dot, on every seed and across the whole range of confinement, to within
$0.35\%$ of the reference or, where there is none, of our own energies. Where it fails, it fails through its
samples: a fixed proposal collapses in the Wigner regime, to $1.5$ effective samples of
$4096$.

Read as instruments, the trained states resolve the crossover shell by shell. Radial shells
form and stiffen across the whole range, and the six-electron ring approaches its classical
radius to within $1.7\%$. Angular order does not keep pace: $n$-fold order within the shells
stays at its random-angle value down to $\omega=0.1$ and becomes large only at
$\omega=10^{-3}$, where the twelve-electron dot is predominantly a $(3,9)$ molecule and the
twenty-electron dot a $(1,7,12)$ one. Those measurements return to the solver. A proposal
built from the measured shell geometry raises the effective sample size fifteen- to
thirtyfold over the fixed one, and carries a Markov-chain-free cascade of converged
six-electron states to $\omega=0.0035$, below which a converged state no longer survives the
step to a weaker trap.

Every state is the lowest found at fixed spin projection, the architecture and optimiser
comparisons rest mainly on $N{=}6$, and the benchmark references also guide pretraining.
\clearpage
